\subsection{Cross Products}

Let K be a field, and let G be a group whose identity element we denote by e. Let L be a commutative K-algebra, and let \(\lambda\) be a homomorphism from G to the automorphism group of the K-algebra L. For any g in G, let \(\tau(g)\) be the automorphism of the multiplicative group \(L^{*}\) of L induced by \(\lambda(g)\) .

Let \(\mathcal{E} = (\Gamma, \iota, \pi)\) be a \(\tau\) -extension of the group \(G\) by \(L^*\) . Define a right action of the group \(L^*\) on the set \(L \times \Gamma\) by

\[
(\beta , \gamma). \alpha = (\beta \alpha , \iota (\alpha) ^ {- 1} \gamma)\tag{20}
\]

for \(\alpha \in \mathrm{L}^*\) , \(\beta \in \mathrm{L}\) , and \(\gamma \in \Gamma\) . We denote by E the set of orbits of \(\mathrm{L}^*\) in \(\mathrm{L} \times \Gamma\) and by \([\beta; \gamma]\) the orbit of the pair \((\beta, \gamma)\) . We therefore have, by construction, the relation

\[
[ \beta \alpha ; \gamma ] = [ \beta ; \iota (\alpha) \gamma ]\tag{21}
\]

for \(\alpha\in\mathrm{L}^{*}\) , \(\beta\in\mathrm{L}\) , and \(\gamma\in\Gamma\) .

Given \(\beta\) in L and \(\gamma\) in \(\Gamma\) , we denote by \(\gamma\beta\) the transform of \(\beta\) under the automorphism \(\lambda\circ\pi(\gamma)\) of L; we have the relations

\[
{ } ^ { \gamma } ( { } ^ { \gamma } { } ^ { \prime } \beta ) = { } ^ { \gamma } { } ^ { \gamma } { } ^ { \prime } \beta   , \quad { } ^ { \gamma } ( \beta + \beta ^ { \prime } ) = { } ^ { \gamma } \beta + { } ^ { \gamma } \beta ^ { \prime }   , \quad \text {and} \quad { } ^ { \gamma } ( \beta \beta ^ { \prime } ) = { } ^ { \gamma } \beta   { } ^ { \gamma } \beta ^ { \prime }\tag{22}
\]

for \(\gamma,\gamma'\) in \(\Gamma\) and \(\beta,\beta'\) in L. There exists a law of composition on E characterized by the relation

\[
[ \beta ; \gamma ] [ \beta^ {\prime}; \gamma^ {\prime} ] = [ \beta^ {\gamma} \beta^ {\prime}; \gamma \gamma^ {\prime} ].\tag{23}
\]

Indeed, it suffices to verify that the right-hand side does not change if we replace, respectively, \(\beta\) , \(\gamma\) , \(\beta'\) , and \(\gamma'\) with \(\beta\alpha\) , \(\iota(\alpha)^{-1}\gamma\) , \(\beta'\alpha'\) , and \(\iota(\alpha')^{-1}\gamma'\) for all \(\alpha\) , \(\alpha'\) in \(L^*\) . This follows immediately by applying to \(\mathcal{E}\) formula (1) of VIII, p.~285, which can also be written as \(\gamma\iota(\alpha) = \iota(\gamma\alpha)\gamma\) for \(\gamma \in \Gamma\) and \(\alpha \in L^*\) . By the formulas in (22), the set E endowed with this law is a monoid with identity element {[}1; e{]}.

Since \(\pi \circ \iota\) is constant with value \(e\) , there exists a mapping \(\widetilde{\pi}\) from E to G such that we have

\[
\widetilde {\pi} ([ \beta ; \gamma ]) = \pi (\gamma)\tag{24}
\]

for \(\beta \in \mathrm{L}\) and \(\gamma \in \Gamma\) .

Let \(g\) be an element of \(\mathrm{G}\) , and let \(\mathrm{E}_g = \widetilde{\pi}^{-1}(g)\) . If \(\gamma_0\) is a fixed element of \(\pi^{-1}(g)\) , then the mapping \(\beta \mapsto [\beta; \gamma_0]\) is a bijection from \(\mathrm{L}\) to \(\mathrm{E}_g\) , thanks to which we will transfer to \(\mathrm{E}_g\) the K-vector space structure on \(\mathrm{L}\) . Since \(\pi^{-1}(g)\)

is composed of elements of the form \(\iota(\alpha)\gamma_{0}\) , where \(\alpha\) runs through \(L^{*}\) , and we have

\[
[ \beta ; \iota (\alpha) \gamma_ {0} ] = [ \beta \alpha ; \gamma_ {0} ],
\]

the vector space structure on \(E_{g}\) does not depend on the choice of \(\gamma_{0}\) .

Let \(g\) and \(g'\) be elements of \(\mathrm{G}\) ; by formulas (22) and (23), the law of composition on \(\mathrm{E}\) induces by restriction a K-bilinear mapping from \(\mathrm{E}_g \times \mathrm{E}_{g'}\) to \(\mathrm{E}_{gg'}\) . Consequently, the vector space \(\mathrm{P} = \bigoplus_{g \in \mathrm{G}} \mathrm{E}_g\) is endowed with the structure of an associative and unital algebra, whose multiplication induces the previous bilinear mapping from \(\mathrm{E}_g \times \mathrm{E}_{g'}\) to \(\mathrm{E}_{gg'}\) for all \(g\) and \(g'\) in \(\mathrm{G}\) . The algebra \(\mathrm{P}\) is called the cross product of \(\mathrm{L}\) and \(\mathcal{E}\) and is denoted by \(\mathbf{A}[\mathcal{E};\mathrm{L}]\) ; its unit element is the element \([1;e]\) of \(\mathrm{E}_e\) .

Set

\[
u (\beta) = [ \beta ; e ]\tag{25}
\]

for \(\beta\) in L. Then \(u: \mathrm{L} \to \mathbf{A}[\mathcal{E};\mathrm{L}]\) is an injective homomorphism of K-algebras. By (23), for every \(\gamma \in \Gamma\) , the element \([1;\gamma]\) is invertible in \(\mathbf{A}[\mathcal{E},\mathrm{L}]\) , and the mapping \(v: \Gamma \to \mathbf{A}[\mathcal{E},\mathrm{L}]^*\) that sends \(\gamma\) to \([1;\gamma]\) is an injective group homomorphism. The homomorphisms \(u\) and \(v\) are called canonical. We have the relations

\begin{enumerate}
\def\labelenumi{(\arabic{enumi})}
\setcounter{enumi}{25}
\tightlist
\item
\end{enumerate}

\[
u (\alpha) = v (\iota (\alpha)),\tag{27}
\]

\[
u (\gamma \beta) = v (\gamma) u (\beta) v (\gamma) ^ {- 1},\tag{28}
\]

\[
[ \beta ; \gamma ] = u (\beta) v (\gamma)
\]

for \(\alpha \in \mathrm{L}^*\) , \(\beta \in \mathrm{L}\) , and \(\gamma \in \Gamma\) .

Conversely, we have the following universal property of the algebra \(A[E;L]\) .

PROPOSITION 12. --- Let B be a K-algebra, \(u' : L \to B\) a K-algebra homomorphism, and \(v' : \Gamma \to B^{*}\) a group homomorphism. We assume satisfied the relations

\begin{enumerate}
\def\labelenumi{(\arabic{enumi})}
\setcounter{enumi}{28}
\tightlist
\item
\end{enumerate}

\[
u ^ {\prime} (\alpha) = v ^ {\prime} (\iota (\alpha)),\tag{30}
\]

\[
u ^ {\prime} (\gamma \beta) = v ^ {\prime} (\gamma) u ^ {\prime} (\beta) v ^ {\prime} (\gamma) ^ {- 1}
\]

for \(\alpha \in \mathrm{L}^*\) , \(\beta \in \mathrm{L}\) , and \(\gamma \in \Gamma\) . Then there exists a unique algebra homomorphism \(f\) from \(\mathbf{A}[\mathcal{E};\mathrm{L}]\) to \(\mathrm{B}\) such that we have \(u' = f \circ u\) and \(v' = f \circ v\) .

To prove the uniqueness of \(f\) , observe that the vector space \(\mathbf{A}[\mathcal{E};\mathbf{L}]\) over the field \(\mathrm{K}\) is generated by the set of elements of the form \([\beta ;\gamma ] = u(\beta)v(\gamma)\) .

If the homomorphism \(f' : A[\mathcal{E}; L] \to B\) satisfies \(f' \circ u = u'\) and \(f' \circ v = v'\) , then it sends \([\beta; \gamma]\) to \(u'(\beta)v'(\gamma)\) and therefore coincides with f.

By formula (29), we have

\[
u ^ {\prime} (\beta \alpha) v ^ {\prime} (\iota (\alpha) ^ {- 1} \gamma) = u ^ {\prime} (\beta) v ^ {\prime} (\gamma)\tag{31}
\]

for \(\alpha\) in \(\mathrm{L}^*\) , \(\beta\) in \(\mathrm{L}\) , and \(\gamma\) in \(\Gamma\) . By the definition of \(\mathrm{E}\) , there consequently exists a mapping \(f_0: \mathrm{E} \to \mathrm{B}\) such that \(f_0([\beta; \gamma]) = u'(\beta)v'(\gamma)\) . By formulas (23) and (30), we have \(f_0(xx') = f_0(x)f_0(x')\) for \(x\) and \(x'\) in \(\mathrm{E}\) . The restriction of \(f_0\) to \(\mathrm{E}_g\) is K-linear for every element \(g\) of \(\mathrm{G}\) ; consequently, there exists a unique K-linear mapping \(f\) from \(\mathbf{A}[\mathcal{E}; \mathrm{L}]\) to \(\mathrm{B}\) that coincides with \(f_0\) on \(\mathrm{E}\) . The mapping \(f\) is an algebra morphism that satisfies \(u' = f \circ u\) and \(v' = f \circ v\) .

Remark. --- Let \(\sigma: G \to \Gamma\) be a section of the mapping \(\pi\) . Set \(\varepsilon_{g} = v(\sigma(g))\) for every g in G, and denote by \(c_{\sigma}\) the 2-cocycle associated with \(\sigma\) (VIII, p.~296). In particular, we have

\[
\varepsilon_ {g} \varepsilon_ {g ^ {\prime}} = u (c _ {\sigma} (g, g ^ {\prime})) \varepsilon_ {g g ^ {\prime}}\tag{32}
\]

for all \(g, g' \in G\) . Let us, moreover, identify L with a subalgebra of \(A[E; L]\) via the homomorphism u. Then every element of \(A[E; L]\) can be written uniquely as \(\sum_{g \in G} a_g \varepsilon_g\) , where \((a_g)\) is a family of elements of L with finite support. The multiplication in \(A[E; L]\) can be expressed by the formula

\[
\left(\sum_ {g} a _ {g} \varepsilon_ {g}\right) \left(\sum_ {g} b _ {g} \varepsilon_ {g}\right) = \sum_ {g} d _ {g} \varepsilon_ {g}\tag{33}
\]

with

\[
d _ {g} = \sum_ {h h ^ {\prime} = g} a _ {h} (\lambda (h) \cdot b _ {h ^ {\prime}}) c _ {\sigma} (h, h ^ {\prime}).\tag{34}
\]

If the extension \(\mathcal{E}\) is semitrivial, then we can choose a section \(\sigma\) of \(\pi\) that is a group morphism from G to \(\Gamma\) ; the cocycle \(c_{\sigma}\) is therefore constant with value 1, and formula (34) simplifies to

\[
d _ {g} = \sum_ {h h ^ {\prime} = g} a _ {h} \lambda (h) \cdot b _ {h ^ {\prime}}.\tag{35}
\]

Let \(K'\) be an extension of the field K. Denote by \(L'\) the \(K'\) -algebra \(L_{(K')}\) , and for any g in G, denote by \(\lambda'(g)\) the automorphism \(\lambda(g)_{(K')}\) induced by \(\lambda(g)\) on \(L_{(K')}\) . Moreover, let us denote by \(\tau'(g)\) the automorphism of \(L'^*\) induced by \(\lambda'(g)\) . Finally, let h be the homomorphism from \(L^*\) to \(L'^*\) that sends x to \(x \otimes 1\) . Let \(\mathcal{E}' = (\Gamma', \iota', \pi')\) be the direct image of the extension of E by h (VIII, p.~289). Let \(A[\mathcal{E}', L']\) be the cross product \(K'\) -algebra of \(\mathcal{E}'\) and \(L'\) , and let \(u': L' \to A[\mathcal{E}', L']\) and \(v': \Gamma' \to A[\mathcal{E}', L']^*\) be the canonical homomorphisms.

PROPOSITION 13. --- There exists a unique isomorphism of \(\mathrm{K}'\) -algebras \(\varphi\) from \(\mathbf{A}[\mathcal{E};\mathrm{L}]_{(\mathrm{K}')}\) to \(\mathbf{A}[\mathcal{E}';\mathrm{L}']\) that satisfies the relations

\[
u ^ {\prime} (h (\beta)) = \varphi (1 \otimes u (\beta))\tag{36}
\]

for \(\beta \in \mathrm{L}\) and

\[
v ^ {\prime} (h (\gamma)) = \varphi (1 \otimes v (\gamma))\tag{37}
\]

for \(\gamma\in\Gamma\) .

The proof follows immediately from the constructions.

The isomorphism \(\varphi\) is called canonical.
% semantic-fragments: legacy-input-seam
\relax{}
